% === ШАБЛОН: ОЦЕНКА РИСКОВ ===
% ГОСТ Р 57580-2017, ISO 27005
\documentclass[14pt,a4paper]{extarticle}

% Подключение базовых настроек
% === ОБЩИЕ НАСТРОЙКИ ДЛЯ ВСЕХ ШАБЛОНОВ ===

\newcommand{\docTitle}{Документ}
\newcommand{\docAuthor}{Организация}
\newcommand{\docHeader}{Заголовок документа}
\newcommand{\docObject}{}
\newcommand{\docCategory}{}
\newcommand{\docCity}{г. Омск}
\newcommand{\docDate}{\today}
\newcommand{\docOrganization}{Организация}
\newcommand{\docType}{Документ}
\newcommand{\docSubtitle}{}
\newcommand{\docSubject}{Тема документа}
\newcommand{\docKeywords}{}
\newcommand{\docYear}{\the\year}

% 2. Шрифты и языки (XeLaTeX)
\usepackage{fontspec}
\usepackage{polyglossia}
\setmainlanguage{russian}
\setotherlanguage{english}

\defaultfontfeatures{Ligatures=TeX, Script=Cyrillic, Path=C:/Windows/Fonts/}
\setmainfont{times.ttf}[
  UprightFont = *,
  BoldFont = timesbd.ttf,
  ItalicFont = timesi.ttf,
  BoldItalicFont = timesbi.ttf
]
\setsansfont{arial.ttf}[
  UprightFont = *,
  BoldFont = arialbd.ttf,
  ItalicFont = ariali.ttf,
  BoldItalicFont = arialbi.ttf
]
\setmonofont{cour.ttf}[
  UprightFont = *,
  BoldFont = courbd.ttf,
  ItalicFont = couri.ttf,
  BoldItalicFont = courbi.ttf
]

% 3. Основные пакеты
\usepackage{amsmath, amssymb, amsfonts}
\usepackage{graphicx}
\usepackage{float}
\usepackage{array}
\usepackage{tabularx}
\usepackage{longtable}
\usepackage{booktabs}
\usepackage{multirow}

% 4. Списки
\usepackage{enumitem}
\setlist[itemize]{label={---}, leftmargin=*, topsep=4pt, itemsep=2pt}
\setlist[enumerate]{label={\arabic*.}, leftmargin=*, topsep=4pt, itemsep=2pt}

% 5. Ссылки и PDF-метаданные (теперь \docTitle уже определён!)
\usepackage{hyperref}
\hypersetup{
    colorlinks=false,
    pdfborder={0 0 0},
    unicode=true,
    pdftitle={\docTitle},
    pdfauthor={\docAuthor},
    pdfsubject={\docSubject}
}

% 6. Совместимость с Pandoc
\providecommand{\tightlist}{%
  \setlength{\itemsep}{0pt}\setlength{\parskip}{0pt}}

% 7. Абзацы
\setlength{\parindent}{1.25cm}
\setlength{\parskip}{0pt}

% 8. Команды для ссылок на нормативку
\newcommand{\gost}[1]{ГОСТ\,#1}
\newcommand{\fstek}[1]{ФСТЭК\,#1}
\newcommand{\fz}[1]{ФЗ\,#1}

% 9. Термины
\newcommand{\term}[1]{\textbf{#1}}
\newcommand{\abbr}[2]{#1 (#2)}

% 10. Стиль для основной страницы
\newcommand{\mainpagestyle}{%
    \pagestyle{fancy}%
    \fancyhf{}%
    \fancyhead[C]{\small \docHeader}%
    \fancyfoot[C]{\thepage}%
    \renewcommand{\headrulewidth}{0.5pt}%
    \renewcommand{\footrulewidth}{0pt}%
}
% === ГЕОМЕТРИЯ ПО ГОСТ Р 7.0.97-2016 ===
% Подключается после preamble.tex

\usepackage{geometry}
\geometry{
    a4paper,
    left=30mm,
    right=10mm,
    top=20mm,
    bottom=20mm,
    headheight=15pt,
    headsep=10mm,
    footskip=10mm,
    includehead,
    includefoot
}

% === Интервалы и шрифт ===
\usepackage{setspace}
\onehalfspacing

% Размер шрифта 14pt с межстрочным интервалом 18pt
\AtBeginDocument{%
  \fontsize{14pt}{18pt}\selectfont
}
% === КОЛОНТИТУЛЫ И ЗАГОЛОВКИ ===
% Подключается после geometry.tex

\usepackage{fancyhdr}

% Базовый стиль колонтитулов
\pagestyle{fancy}
\fancyhf{}
\fancyhead[C]{\small \docHeader}
\fancyfoot[C]{\thepage}
\renewcommand{\headrulewidth}{0.5pt}
\renewcommand{\footrulewidth}{0pt}

% === Команды для заголовков ===
\newcommand{\makeDocHeader}{%
    \begin{center}
        \Large\bfseries \docTitle \\
        \normalsize \docSubtitle \\
        \vspace{0.5cm}
        \small \docAuthor \\
        \docDate
    \end{center}
    \vspace{1cm}
}

\newcommand{\docSection}[1]{%
    \section{#1}%
    \addcontentsline{toc}{section}{\thesection\quad #1}%
}

\newcommand{\docSubsection}[1]{%
    \subsection{#1}%
    \addcontentsline{toc}{subsection}{\thesubsection\quad #1}%
}

% === Специальные стили страниц ===
\fancypagestyle{titlepage}{%
    \fancyhf{}%
    \renewcommand{\headrulewidth}{0pt}%
    \renewcommand{\footrulewidth}{0pt}%
}

\fancypagestyle{tocpage}{%
    \fancyhf{}%
    \fancyhead[C]{\small \docHeader}%
    \fancyfoot[C]{\thepage}%
    \renewcommand{\headrulewidth}{0.5pt}%
}
% === ТИТУЛЬНЫЙ ЛИСТ ===
% Подключается после headers.tex

% Основной титульный лист (для политик, моделей угроз, оценки рисков)
\newcommand{\makeTitlePage}{%
\begin{titlepage}
\pagestyle{empty}
\centering
\vfill
{\normalsize\bfseries \docOrganization\par}
\vspace{2cm}
{\Large\bfseries \docType\par}
\vspace{0.5cm}
{\Huge\bfseries \docTitle\par}
\vspace{1cm}

\ifx\docObject\empty\else
{\normalsize Объект защиты: \docObject\par}
\vspace{0.3cm}
\fi

\ifx\docCategory\empty\else
{\normalsize Категория: \docCategory\par}
\vspace{0.3cm}
\fi

\vfill
\begin{flushright}
\normalsize
СОГЛАСОВАНО:\\
Руководитель службы ИБ\\
\vspace{1.5cm}
\underline{\hspace{6cm}}\\
\small «\underline{\hspace{1.5cm}}» \underline{\hspace{3cm}} 20\underline{\hspace{1cm}} г.
\end{flushright}
\vspace{2cm}
{\normalsize \docCity, \docYear\par}
\vspace{1cm}
\vfill
\end{titlepage}
}

% Упрощённый титульный лист (для инструкций, регламентов)
\newcommand{\makeSimpleTitlePage}{%
\begin{titlepage}
\pagestyle{empty}
\centering
\vfill
{\normalsize\bfseries \docOrganization\par}
\vspace{2cm}
{\Large\bfseries \docType\par}
\vspace{0.5cm}
{\Huge\bfseries \docTitle\par}
\vspace{1cm}

\ifx\docObject\empty\else
{\normalsize Объект: \docObject\par}
\vspace{0.3cm}
\fi

\vfill
\begin{flushright}
\normalsize
УТВЕРЖДАЮ:\\
Руководитель организации\\
\vspace{1.5cm}
\underline{\hspace{6cm}}\\
\small «\underline{\hspace{1.5cm}}» \underline{\hspace{3cm}} 20\underline{\hspace{1cm}} г.
\end{flushright}
\vspace{2cm}
{\normalsize \docCity, \docYear\par}
\vspace{1cm}
\vfill
\end{titlepage}
}

% Переменные Pandoc
$if(title)$\renewcommand{\docTitle}{$title$}$endif$
$if(author)$\renewcommand{\docAuthor}{$author$}$endif$
$if(doc_header)$\renewcommand{\docHeader}{$doc_header$}$endif$
$if(date)$\renewcommand{\docDate}{$date$}$endif$
$if(city)$\renewcommand{\docCity}{$city$}$endif$
$if(year)$\renewcommand{\docYear}{$year$}$endif$

% Специфичные переменные
\newcommand{\docType}{ОЦЕНКА РИСКОВ}
\newcommand{\docSubtitle}{информационной безопасности}
\newcommand{\docSubject}{Анализ и оценка рисков}
\newcommand{\docKeywords}{риски, оценка, обработка рисков}

\begin{document}

% Титульный лист
\makeTitlePage

% Оглавление
\mainpagestyle
\tableofcontents
\newpage

% 1. Введение
\section{Введение}
\label{sec:intro}

$body$

% 2. Методология
\section{Методология оценки рисков}
\label{sec:methodology}

$if(methodology)$
$methodology$
$else$
Оценка рисков проводится в соответствии с:
\begin{itemize}
    \item ГОСТ Р 57580-2017;
    \item ISO/IEC 27005:2018;
    \item Методикой ФСТЭК России.
\end{itemize}

\subsection{Подход к оценке}
Применяется комбинированный подход:
\begin{itemize}
    \item \textbf{Качественная оценка} — экспертный анализ;
    \item \textbf{Количественная оценка} — расчет показателей.
\end{itemize}
$endif$

% 3. Идентификация активов
\section{Идентификация активов}
\label{sec:assets}

$if(assets)$
$assets$
$else$
\subsection{Критические активы}
\begin{table}[h]
\centering
\caption{Реестр активов}
\begin{tabular}{|l|l|l|}
\hline
\textbf{Актив} & \textbf{Тип} & \textbf{Критичность} \\
\hline
База данных ПДн & Информация & Высокая \\
Серверы & Оборудование & Высокая \\
Сеть & Инфраструктура & Средняя \\
\hline
\end{tabular}
\end{table}
$endif$

% 4. Идентификация угроз
\section{Идентификация угроз}
\label{sec:threats}

$if(threats)$
$threats$
$endif$

% 5. Идентификация уязвимостей
\section{Идентификация уязвимостей}
\label{sec:vulnerabilities}

$if(vulnerabilities)$
$vulnerabilities$
$endif$

% 6. Анализ рисков
\section{Анализ рисков}
\label{sec:analysis}

$if(analysis)$
$analysis$
$else$
\subsection{Матрица рисков}
\begin{table}[h]
\centering
\caption{Оценка рисков}
\begin{tabular}{|l|l|l|l|}
\hline
\textbf{Риск} & \textbf{Вероятность} & \textbf{Влияние} & \textbf{Уровень} \\
\hline
Утечка ПДн & Средняя & Высокое & Высокий \\
DDoS-атака & Низкая & Среднее & Средний \\
Отказ оборудования & Средняя & Высокое & Высокий \\
\hline
\end{tabular}
\end{table}
$endif$

% 7. Оценка рисков
\section{Оценка рисков}
\label{sec:risk_evaluation}

$if(risk_evaluation)$
$risk_evaluation$
$endif$

% 8. Реестр рисков
\section{Реестр рисков}
\label{sec:risk_register}

$if(risk_register)$
$risk_register$
$else$
\begin{table}[h]
\centering
\caption{Реестр выявленных рисков}
\begin{tabular}{|c|p{6cm}|c|c|}
\hline
\textbf{№} & \textbf{Описание риска} & \textbf{Уровень} & \textbf{Приоритет} \\
\hline
1 & Несанкционированный доступ к ПДн & Высокий & 1 \\
2 & Потеря данных из-за сбоя & Высокий & 1 \\
3 & Инсайдерская угроза & Средний & 2 \\
\hline
\end{tabular}
\end{table}
$endif$

% 9. Обработка рисков
\section{План обработки рисков}
\label{sec:treatment}

$if(treatment)$
$treatment$
$else$
\subsection{Стратегии обработки}
Для каждого риска выбирается одна из стратегий:
\begin{itemize}
    \item \textbf{Снижение} — внедрение мер защиты;
    \item \textbf{Принятие} — осознанное принятие риска;
    \item \textbf{Передача} — страхование, аутсорсинг;
    \item \textbf{Избегание} — отказ от рискованной деятельности.
\end{itemize}
$endif$

% 10. Меры защиты
\section{Меры по снижению рисков}
\label{sec:controls}

$if(controls)$
$controls$
$endif$

% Приложения
$if(applications)$
\newpage
\appendix
\section*{Приложения}
\addcontentsline{toc}{section}{Приложения}
$applications$
$endif$

% Нормативные ссылки
$if(standards)$
\newpage
\section*{Нормативные ссылки}
\addcontentsline{toc}{section}{Нормативные ссылки}
\begin{itemize}
$for(standards)$\item $standards$$endfor$
\end{itemize}
$endif$

\end{document}